\documentclass[12pt]{article}
\usepackage[utf8]{inputenc}
\usepackage[T1]{fontenc}
\usepackage[margin=1in]{geometry}
\usepackage{setspace}
\usepackage{xcolor}
\usepackage[colorlinks=true, urlcolor=blue!60!black, linkcolor=blue!60!black]{hyperref}
\setlength{\parindent}{2em}
\setlength{\parskip}{1.25em}
\renewcommand{\baselinestretch}{1.0}
\begin{document}
\begin{center}
{\LARGE \textit{Williamson v. Lee Optical of Oklahoma, Inc.}, 348 U.S. 483 (1955)\par}
\vspace{0.8em}
{\large Case Brief --- Four-Page Edition\par}
\end{center}
\vspace{1em}

\section{Memory Jogger}

Under the Fourteenth Amendment's Due Process and Equal Protection Clauses, economic regulations are upheld under hyper-deferential rational basis review whenever a court can conceive of any plausible legislative rationale---even if the law seems needless, wasteful, or addresses a problem only one step at a time (\href{https://www.courtlistener.com/opinion/105285/williamson-v-lee-optical-of-oklahoma-inc/}{348 U.S. 483}).

\section{Detailed Case Facts}

In 1953 Oklahoma enacted a comprehensive statute regulating visual care and the optical trade. The industry has three roles: ophthalmologists (physicians specializing in eye care), optometrists (who examine eyes for refractive error and fill prescriptions), and dispensing opticians (artisans who grind lenses and fit frames). Four provisions were challenged:
\begin{enumerate}
\item \textbf{Prescription requirement (\S~2):} only a licensed ophthalmologist or optometrist could fit lenses or frames or duplicate or replace lenses, except on a written prescription; an optician could not put old lenses into new frames or duplicate a broken lens without one.
\item \textbf{Ready-to-wear exemption (\S~3):} sellers of ready-to-wear glasses, plano sunglasses, and safety goggles were exempt.
\item \textbf{Advertising ban (\S~3):} soliciting the sale of eyeglasses, lenses, frames, or eye examinations through radio, television, window displays, and other media was forbidden.
\item \textbf{Retail space ban (\S~4):} retail merchandisers could not rent space to anyone purporting to examine eyes or provide visual care.
\end{enumerate}
Lee Optical and other opticians sued Attorney General Mac Q. Williamson and other officials for a declaratory judgment and injunction under the Due Process and Equal Protection Clauses.

\section{Procedural History}

A three-judge District Court invalidated the prescription requirement for fitting old lenses and duplicating broken ones, finding the means ``neither reasonably necessary nor reasonably related to the end sought to be achieved''; struck down the ready-to-wear exemption under equal protection; and invalidated the frame-advertising ban (over Judge Murrah's dissent) and the retail space ban. It upheld the ban on advertising spectacles, lenses, and prisms. Both sides appealed directly; the Supreme Court heard argument on March 2, 1955, and decided the case on March 28.

\section{Judicial Votes}

\textbf{8--0}, with one Justice not participating.
\begin{itemize}
\item \textbf{Majority:} Justice Douglas, joined by Chief Justice Warren and Justices Black, Reed, Frankfurter, Burton, Clark, and Minton.
\item \textbf{Concurrences and dissents:} None.
\item \textbf{Not participating:} Justice Harlan.
\end{itemize}

\section{Holding}

Affirmed in part and reversed in part: every challenged provision is constitutional. A state economic regulation satisfies due process if there is an evil at hand for correction and the legislature might rationally have thought its measure a reasonable way to correct it, however unwise or imperfect. An economic classification violates equal protection only if it amounts to ``invidious discrimination,'' and a legislature may reform one phase of a field at a time. Douglas's declaration that ``[t]he day is gone'' for substantive due process invalidation of business regulation, and his invocation of \textit{Munn v. Illinois}, are sweeping dicta signaling the burial of \textit{Lochner}, though they did not foreclose later commercial speech challenges.

\section{Analysis of Opinions}

\textbf{Majority (Douglas): prescription requirement.} Relying on \textit{Roschen v. Ward}, Douglas conceded that the requirement ``may exact a needless, wasteful requirement in many cases,'' but weighing costs and benefits is for the legislature. It \textit{might} have concluded that prescriptions are needed often enough to justify a blanket rule, or that every frame change should prompt a visit to a specialist who could detect latent disease. That the law did not actually require a new examination did not matter, since a law ``need not be in every respect logically consistent with its aims.'' Citing \textit{Nebbia} and \textit{West Coast Hotel}, he declared the era of \textit{Lochner}-style invalidation over.

\textbf{Ready-to-wear exemption.} Classification admits of ``no doctrinaire definition''; reform may proceed ``one step at a time,'' and ready-to-wear glasses might not loom large or might pose distinct problems.

\textbf{Advertising and retail space.} Adopting Judge Murrah's dissent, Douglas reasoned that frames are used with lenses, so frame advertisements might bait customers into buying lenses; under \textit{Semler}, all who deal with the eye may be treated as professionals. Barring eye-care providers from retail stores could free the profession from ``all taints of commercialism.'' The lens-advertising ban was upheld for the same reasons.

\textbf{Separate opinions.} None were filed.

\section{Examples: Future Applications}

\begin{enumerate}
\item \textbf{Same-Side (Upheld): Casket-Seller Licensing.} Requiring casket retailers to obtain a full funeral director's license, while exempting out-of-state online sellers, survives if a court can hypothesize a consumer-protection rationale, even if it looks like protectionism.
\item \textbf{Same-Side (Upheld): Attendant-Only Gas Stations.} Banning self-service gasoline while letting electric-vehicle drivers charge themselves is rational, since fuel poses distinct hazards and reform may proceed one step at a time.
\item \textbf{Opposite-Side (Invalid): Racial Licensing Bar.} Barring a racial or national-origin group from optician licenses triggers strict scrutiny; \textit{Lee Optical} covers only non-invidious economic classifications.
\item \textbf{Opposite-Side (Outside \textit{Lee Optical}): Price-List Ban.} A ban on truthful online price lists, challenged under the First Amendment, faces commercial speech scrutiny rather than rational-basis deference.
\item \textbf{Fence-Sitter: Hair-Braiding Licensure.} Requiring braiders to complete 1,500 hours of cosmetology school with no evidence of health risk could be upheld on an imagined sanitation rationale, but some courts applying ``rational basis with bite'' question whether naked protectionism survives.
\end{enumerate}

\section{Critique}

\textbf{Libertarian and originalist critiques.} Randy Barnett and advocates at the Cato Institute, Pacific Legal Foundation, and Institute for Justice argue that \textit{Lee Optical} abdicated the duty to protect the right to earn an honest living from ``class legislation,'' turning rational basis review into a rubber stamp for special-interest cartels by inventing rationales the legislature never had. A 2025 \textit{Harvard Law Review} Note argues the Court abandoned the more balanced \textit{Carolene Products}-era standard, which still policed naked protectionism.

\textbf{Progressive and restraint views.} Defenders see the decision as essential to preventing a revival of \textit{Lochner}, leaving economic tradeoffs to elected legislatures. Some progressives note the irony that the same deference leaves low-income entrepreneurs exposed to captured political processes.

\section{Key Quotations}

\begin{quote}
``But the law need not be in every respect logically consistent with its aims to be constitutional. It is enough that there is an evil at hand for correction, and that it might be thought that the particular legislative measure was a rational way to correct it.''
\end{quote}
\noindent --- Justice Douglas, Opinion of the Court, 348 U.S. at 487--488.

\begin{quote}
``The problem of legislative classification is a perennial one, admitting of no doctrinaire definition. Evils in the same field may be of different dimensions and proportions, requiring different remedies. Or so the legislature may think.''
\end{quote}
\noindent --- Justice Douglas, Opinion of the Court, 348 U.S. at 489.

\begin{quote}
``The day is gone when this Court uses the Due Process Clause of the Fourteenth Amendment to strike down state laws, regulatory of business and industrial conditions because they may be unwise, improvident, or out of harmony with a particular school of thought.''
\end{quote}
\noindent --- Justice Douglas, Opinion of the Court, 348 U.S. at 488.

\end{document}
